% Source: https://artofproblemsolving.com/community/c3h1978070_factorization_for_sum_of_squares
% Stable identity: aops-topic-1978070-q01
(a) When $ab\ge0$, $\sqrt{2ab}$ is real and by the difference of squares
\[
\left(a+b+\sqrt{2ab}\right)\left(a+b-\sqrt{2ab}\right)=(a+b)^2-2ab=a^2+b^2 .
\]
When $ab<0$, $\sqrt{2ab}$ is not a real number. So the equation holds, with real quantities, exactly when $\boxed{ab\ge0}$. (If one allows $\sqrt{2ab}=i\sqrt{-2ab}$, the algebra still works, but then the factors are complex.)

(b) The factors involve $\sqrt{2ab}$, which is not a polynomial in $a$ and $b$, so this is an identity between functions on part of the plane, not a polynomial factorization.

Suppose $a^2+b^2=P\cdot Q$ with $P$, $Q$ non-constant real polynomials. Degrees add, so both have degree $1$, and the product of their degree-$1$ parts is the degree-$2$ part $a^2+b^2$:
\[
a^2+b^2=(pa+qb)(sa+tb)\qquad(p,q,s,t\in\mathbb R).
\]
Comparing coefficients of $a^2$ gives $ps=1$, so $p\ne0$. Put $b=1$ and $a=-\dfrac {q}{p}$: the right side is $0$, while the left side is $\dfrac{q^2}{p^2}+1>0$. This contradiction shows that $a^2+b^2$ has no factorization into non-constant real polynomials. $\blacksquare$
